\section{Method}

\subsection{Overview}
\label{sec:method:overview}

Given a long-horizon task $\mathcal{T}$ and a computer environment,
LongHorizon-Harness executes the task through a sequence of dynamically
determined rounds rather than a single continuously growing session.
The harness maintains an explicit task state outside task execution and
advances it only with evidence independently verified from the
environment. Across rounds, only the task state and its supporting audit
reports persist; the executor's raw interaction trajectory is
discarded after each round. Fig.~\ref{fig:pipeline} provides an overview
of the framework.

Each round follows a \emph{Manage-Execute-Audit} (MEA) loop. Let $S_i$
denote the task state available at the beginning of round $i$,
$e_{i-1}\in\mathcal{E}$ the current environment state, and
$V_{i-1}=(v_1,\ldots,v_{i-1})$ the accumulated audit reports. The
manager constructs a bounded subtask contract $c_i$. A fresh-context
executor performs the contract, transforms the environment from
$e_{i-1}$ to $e_i$, and returns an execution report $o_i$. A separate
auditor then inspects $e_i$ through read-only tools and produces audit
report $v_i$. The manager incorporates $v_i$ into the next task state
$S_{i+1}$ before determining whether another round is required. The loop
ends when the audited state satisfies the original task, no
permitted subtask can advance the remaining requirements, user input is
required, or the round budget is exhausted.


\subsection{Manager}
\label{sec:method:manager}

The manager owns the persistent task state and determines how the task
should proceed. It has access to the original task $\mathcal{T}$, the
current task state, and all accumulated audit reports, but has no direct
interface to the computer environment. It cannot observe application
state, inspect workspace contents, invoke GUI or CLI tools, or modify
the environment. Its decisions are based entirely on the task state and
environment evidence recorded by the auditors.

After round $i$, the manager updates the task state and produces the
next control decision:
\begin{equation}
\label{eq:manager}
(S_{i+1},q_{i+1},c_{i+1})
=
\Phi_{\mathrm{mgr}}
\left(
\mathcal{T},
S_i,
V_i
\right),
\end{equation}
where $V_i=(v_1,\ldots,v_i)$ and $q_{i+1}$ is one of
\textsc{execute}, \textsc{done}, \textsc{blocked}, and \textsc{ask}.
The contract $c_{i+1}$ is returned only when further execution is
required.

\paragraph{Task-state update.}
The task state is a structured collection of task-relevant records:
a \textit{requirement} represents an objective or constraint derived
from the original task, an \textit{artifact} represents an output
created or modified during execution, and a \textit{fact} records
environment information needed by subsequent rounds. Each record is
marked as \textit{completed}, \textit{pending}, \textit{blocked}, or
\textit{untrusted}, and retains references to the audit evidence
supporting its current status. The initial state $S_1$ is constructed
from $\mathcal{T}$ with its requirements marked as pending. After each
round, the manager applies the verified findings in $v_i$ to $S_i$,
adding or updating the corresponding records while leaving unresolved
ones pending, blocked, or untrusted as appropriate. Executor claims do
not directly change the persistent state: a record is marked as
completed only when supported by clean audit evidence. 

\paragraph{Next-subtask construction.}
The manager compares $S_{i+1}$ against the original task, selects an
unresolved objective that can be advanced from the current state, and
checks its dependencies and prerequisites. It then constructs a bounded
contract $c_{i+1}$ specifying the immediate goal, acceptance criteria,
boundary constraints, and the task-state records and prior audit reports
relevant to execution and verification. The contract is routed to a GUI
or CLI executor according to the primary environment transition it
requires. The manager returns \textsc{done} when the audited state
satisfies $\mathcal{T}$ without unresolved integrity violations,
\textsc{blocked} when no permitted action can advance the remaining
requirements, and \textsc{ask} when progress requires user information
or authorization; otherwise, it returns \textsc{execute} together with
$c_{i+1}$.

\subsection{Executor}
\label{sec:method:executor}

The executor performs the contract selected by the manager and is the
only role permitted to intentionally modify the environment. In round
$i$, it receives the original task $\mathcal{T}$, the current task
state $S_i$, the subtask contract $c_i$, and only the prior audit
reports referenced by the contract. It transforms the environment from
$e_{i-1}$ to $e_i$:
\begin{equation}
\label{eq:executor}
(e_i,o_i)
=
\Phi_{\mathrm{exec}}
\left(
\mathcal{T},
S_i,
c_i;
e_{i-1}
\right),
\end{equation}
where $o_i$ summarizes the actions performed, resulting state,
artifacts produced or modified, and issues encountered during
execution. The report $o_i$ describes the executor's outcome but does
not establish that the contract has been completed.

\paragraph{Fresh-context execution.}
Each executor invocation runs as a fresh, budget-bounded episode
containing only the information supplied for the current round. It does
not receive the raw interaction trajectories of earlier rounds. Within
the episode, the executor may perform multiple cycles of planning,
environment interaction, observation, and revision. When the episode
ends, its raw trajectory and internal reasoning are discarded; only
$o_i$ is forwarded for auditing.

\paragraph{GUI and CLI capability boundaries.}
GUI and CLI executors operate through different environment interfaces.
The GUI executor receives screen-oriented capabilities such as taking
screenshots, clicking, scrolling, and entering text, and is responsible
for transitions centered on application and interface state. The CLI
executor receives capabilities such as shell execution, file editing,
coding, and testing, and is responsible for transitions centered on
workspace, process, and program state.

The harness exposes only the environment interface and tool set
assigned to the selected executor role. Capabilities outside that role
are unavailable unless explicitly provided by its configuration. This
separates responsibility for different classes of state-changing
actions while allowing the manager to select the interface appropriate
for each contract.

\paragraph{Backend execution.}
Executors are instantiated through a common agent-adapter interface.
Given a contract, a role-specific environment interface, and an
execution budget, the adapter launches an existing backend such as
Claude Code, Codex CLI, or OpenClaw as one bounded episode. The backend
retains its native planning and tool-use loop and may normally invoke
the shell, edit files, write code, run tests, or interact with
applications when those capabilities are exposed by its assigned role.
The harness does not replace the backend's internal execution process;
it controls the supplied context, available tools, environment
permissions, execution budget, and returned report.


\subsection{Auditor}
\label{sec:method:auditor}

The auditor independently verifies the environment state produced by
the executor. After execution, it receives the original task
$\mathcal{T}$, the task state $S_i$, the contract $c_i$, the prior
audit reports referenced by the contract, and the executor report
$o_i$. It does not receive the executor's raw interaction trajectory or
internal reasoning. The auditor inspects the resulting environment
$e_i$ and produces
\begin{equation}
\label{eq:auditor}
v_i
=
\Phi_{\mathrm{aud}}
\left(
\mathcal{T},
S_i,
c_i,
o_i;
e_i
\right),
\end{equation}
where $v_i$ is appended to the persistent audit history and returned to
the manager before the next round.

\paragraph{Independent environment inspection.}
The auditor starts from a fresh context that excludes the executor's
raw interaction trajectory and internal reasoning. It may use $o_i$ to
locate relevant files, windows, logs, processes, or other outputs, but
determines completion by independently comparing the resulting
environment against the goal, acceptance criteria, and boundary
constraints in $c_i$. A GUI auditor examines application and screen
state through observation-oriented interactions, whereas a CLI auditor
uses non-mutating commands and inspection tools to examine files,
metadata, logs, processes, tests, and workspace state. In both cases,
audit conclusions must be supported by evidence obtained directly from
the environment rather than by the executor's completion claim.

\paragraph{Read-only authority.}
The auditor may change its observation view when necessary for
inspection but cannot modify task-relevant environment state. It cannot
create, edit, overwrite, move, or delete protected artifacts, execute
state-changing commands, or perform GUI actions that alter the result
under inspection. The harness monitors task-relevant workspace and
artifact state throughout auditing; any detected mutation is recorded
as an integrity violation, and the resulting report cannot support a
completed task-state record.

\paragraph{Audit-result construction.}
The audit report $v_i$ records three classes of findings. First, it
assigns a completion status of \textit{complete},
\textit{incomplete}, or \textit{blocked} by evaluating the contract's
acceptance criteria. Second, it assigns an integrity status of
\textit{clean}, \textit{suspect}, or \textit{violation} by checking
workspace mutations, artifact validity and provenance, and relevant
deletion constraints. Third, it records the task-state updates supported
by the inspection, including verified facts, supporting evidence, and
remaining gaps.

The auditor may propose changes to requirement, artifact, and fact
records, but the manager determines how these findings are incorporated
into $S_{i+1}$. Consequently, $o_i$ remains an unverified execution
summary, whereas $v_i$ provides the environment-grounded evidence that
may advance the persistent task state.